\section{The norm of the loss}\label{sec:norm}

This section tests the explanation of \cref{sec:theory} directly. The same transformer, with
the same labelled instances, initial weights, order of the instances and learning rate, is
trained with the relative displacement error $\mathcal{L}_D$ \eqref{eq:ld} or with the
relative energy-norm error $\mathcal{L}_K$ \eqref{eq:lk}; nothing else changes
(\cref{sec:methods:training}). Both comparisons were pre-registered, with directions fixed
in advance. \Cref{fig:normlabels} summarises them, and \cref{tab:3d,tab:2d} give every
metric.

\subsection{Three dimensions}\label{sec:norm:three}

% generated from the committed records; do not edit
\begin{table}[tbp]
\centering
\footnotesize
\setlength{\tabcolsep}{4pt}
\caption{Three-dimensional solids: every metric on the 256 validation instances, four load cases each, for the networks trained on 1,024 instances (the label-free transformer on the same instances without their solutions) and for the baselines. Trained networks: mean $\pm$ sample standard deviation over three seeds. Metrics as in \cref{sec:methods:metrics}; lower is better except in critical-region recall.}
\label{tab:3d}
\resizebox{\ifdim\width>\textwidth\textwidth\else\width\fi}{!}{%
\begin{tabular}{lcccccc}
\toprule
Model & \shortstack{Displacement \\ error} & \shortstack{Relative \\ energy gap} & \shortstack{von Mises \\ error} & \shortstack{Peak stress \\ error} & \shortstack{Critical-region \\ recall} & \shortstack{Worse than \\ zero field} \\
\midrule
Label-free & 0.0300 $\pm$ 0.0028 & 0.00912 $\pm$ 0.00228 & 0.0554 $\pm$ 0.0045 & 0.0807 $\pm$ 0.0134 & 0.793 $\pm$ 0.022 & 0/768 \\
Supervised, displacement & 0.0470 $\pm$ 0.0038 & 31.2 $\pm$ 53.5 & 0.370 $\pm$ 0.140 & 0.784 $\pm$ 0.807 & 0.619 $\pm$ 0.005 & 18/768 \\
Supervised, displacement + energy term & 0.0470 $\pm$ 0.0122 & 0.116 $\pm$ 0.050 & 0.128 $\pm$ 0.008 & 0.115 $\pm$ 0.026 & 0.730 $\pm$ 0.006 & 2/768 \\
Supervised, energy norm & 0.0227 $\pm$ 0.0022 & 0.00504 $\pm$ 0.00022 & 0.0440 $\pm$ 0.0014 & 0.0386 $\pm$ 0.0011 & 0.858 $\pm$ 0.001 & 0/768 \\
Graph network, displacement & 0.0168 $\pm$ 0.0001 & 0.447 $\pm$ 0.084 & 0.296 $\pm$ 0.010 & 0.488 $\pm$ 0.011 & 0.679 $\pm$ 0.009 & 28/768 \\
\midrule
Nearest-neighbour field & 0.382 & 1.22 & 0.717 & 1.23 & 0.509 & 88/256 \\
Scaled polynomial & 1.56 & 4,362 & 14.9 & 54.8 & 0.137 & 256/256 \\
\midrule
Ritz projection, $p = 6$ & 0.0208 & 0.0183 & 0.0854 & 0.321 & 0.811 & 0/256 \\
Ritz projection, $p = 7$ & 0.0134 & 0.0114 & 0.0624 & 0.283 & 0.859 & 0/256 \\
Ritz projection, $p = 8$ & 0.0100 & 0.00838 & 0.0521 & 0.253 & 0.879 & 0/256 \\
Ritz projection, $p = 9$ & 0.00798 & 0.00654 & 0.0449 & 0.223 & 0.896 & 0/256 \\
\midrule
Zero field & 1.00 & 1.00 & 1.00 & 1.00 & -- & -- \\
\bottomrule
\end{tabular}}

\smallskip
\begin{minipage}{\linewidth}\footnotesize\raggedright
The Ritz projection of degree $p$ (\cref{sec:methods:ritz}) was added after the pre-registered runs. It, the nearest-neighbour field and the scaled polynomial are deterministic, so their last column counts instances.
\end{minipage}
\end{table}


The three-dimensional energy-norm run trained the transformer on $\mathcal{L}_K$ with the
1,024 labelled instances of the three-dimensional comparison, and evaluated the stored
label-free and displacement-supervised transformers beside it, which reproduced their
earlier values exactly. Its two hypotheses were supported. The energy-norm-supervised
transformer's median relative energy gap, \numKdKnormGapMedian{} averaged over the seeds,
was \numKdGapRel{} lower than the displacement-supervised transformer's,
\numKdLabGapMedian{}, beyond the noise guard of \numKdGapTau{}, a \numNlThreeNormE{}-fold
reduction of the median energy-norm error; and its median von Mises error was \numKdVmRel{}
lower (guard \numKdVmTau{}), in each seed. Its displacement error fell too, from
\numThreeLabDisp{} to \numThreeKnormDisp{} in the seed mean (\cref{tab:3d}), and none of its
\numThreePairs{} instance-seed predictions was worse than the zero field, against
\numThreeLabWorse{} of the displacement-supervised transformer's. The verdicts read each
seed's median over the 256 validation instances, a statistic fixed after the per-instance
values of the two stored transformers, and the two-dimensional comparison's readings on the
medians, were known. The reason is one validation instance, No.~\numWorstIndex{}: in seed
\numThreeLabWorstSeed{} the displacement-supervised transformer's prediction of it had a
relative energy gap of \numThreeLabWorst{}, averaged over its four load cases, which carries
that network's seed means (\cref{tab:3d}) and widens the guard on them. On the seed means
the von Mises error was \numKdMeansVmRel{} lower, beyond the guard of \numKdMeansVmTau{},
although with three seeds its Welch 95\% interval runs from \numKdMeansVmWelchLow{} to
\numKdMeansVmWelchHigh{}, and the relative energy gap \numKdMeansGapRel{} lower, within a
guard that this instance widens to \numKdMeansGapTau{}. On the \numHoldThreeN{} held-out
instances of \cref{sec:holdout}, read without a verdict, the two reductions were
\numKdHoldGapRel{} and \numKdHoldVmRel{}, beyond guards of \numKdHoldGapTau{} and
\numKdHoldVmTau{}. The test rests on one model size and one family of geometries, with the
displacement-supervised settings unchanged and not tuned for the new loss; it supports the
explanation of \cref{sec:theory} rather than establishing it in general.

\subsection{Two dimensions}\label{sec:norm:two}

% generated from the committed records; do not edit
\begin{table}[tbp]
\centering
\footnotesize
\setlength{\tabcolsep}{4pt}
\caption{Two-dimensional plates: every metric on the 256 validation instances, four load cases each, for the networks trained on 1,024 instances (the label-free transformer on the same instances without their solutions) and for the baselines. Trained networks: mean $\pm$ sample standard deviation over three seeds. Metrics as in \cref{sec:methods:metrics}; lower is better except in critical-region recall.}
\label{tab:2d}
\resizebox{\ifdim\width>\textwidth\textwidth\else\width\fi}{!}{%
\begin{tabular}{lcccccc}
\toprule
Model & \shortstack{Displacement \\ error} & \shortstack{Relative \\ energy gap} & \shortstack{von Mises \\ error} & \shortstack{Peak stress \\ error} & \shortstack{Critical-region \\ recall} & \shortstack{Worse than \\ zero field} \\
\midrule
Label-free & 0.0634 $\pm$ 0.0031 & 0.0383 $\pm$ 0.0019 & 0.123 $\pm$ 0.002 & 0.128 $\pm$ 0.001 & 0.753 $\pm$ 0.001 & 0/768 \\
Supervised, displacement & 0.0505 $\pm$ 0.0023 & 0.133 $\pm$ 0.014 & 0.235 $\pm$ 0.015 & 0.141 $\pm$ 0.007 & 0.668 $\pm$ 0.015 & 0/768 \\
Supervised, energy norm & 0.0547 $\pm$ 0.0094 & 0.0290 $\pm$ 0.0125 & 0.106 $\pm$ 0.020 & 0.0947 $\pm$ 0.0356 & 0.804 $\pm$ 0.055 & 0/768 \\
Graph network, displacement & 0.0499 $\pm$ 0.0003 & 0.825 $\pm$ 0.020 & 0.578 $\pm$ 0.009 & 0.321 $\pm$ 0.011 & 0.517 $\pm$ 0.006 & 187/768 \\
\midrule
Nearest-neighbour field & 0.404 & 0.959 & 0.735 & 0.643 & 0.515 & 94/256 \\
Scaled polynomial & 1.42 & 873 & 11.4 & 30.8 & 0.159 & 256/256 \\
\midrule
Ritz projection, $p = 6$ & 0.0500 & 0.0404 & 0.122 & 0.272 & 0.828 & 0/256 \\
Ritz projection, $p = 8$ & 0.0424 & 0.0333 & 0.109 & 0.208 & 0.852 & 0/256 \\
\midrule
Zero field & 1.00 & 1.00 & 1.00 & 1.00 & -- & -- \\
\bottomrule
\end{tabular}}

\smallskip
\begin{minipage}{\linewidth}\footnotesize\raggedright
The Ritz projection of degree $p$ (\cref{sec:methods:ritz}) was added after the pre-registered runs. It, the nearest-neighbour field and the scaled polynomial are deterministic, so their last column counts instances.
\end{minipage}
\end{table}


The two-dimensional comparison trained the displacement-supervised and the
energy-norm-supervised transformers, which differ only in the loss, and pre-registered two
hypotheses about them: H2a, that the energy-norm-supervised transformer has the lower
relative energy gap, which is close to a check that the change of loss took effect, and H2b,
that it has the lower von Mises error. The von Mises error is a term of neither loss: H2b
tests the explanation that \cref{prop:stress} gives. The comparison was evaluated on the
validation instances on which the results that motivated it had been seen, and was
registered with the values of its label-free transformers on them known and after a
functional check at a far smaller scale had shown the directions of its hypotheses
(supplementary Section~\SuppPrereg{}). Both were supported on the validation instances. With
1,024 labelled instances the energy-norm-supervised transformer's relative energy gap was
\numHTwoARel{} lower (guard \numHTwoATau{}), \numHTwoARootRel{} lower in the
root-mean-square relative energy-norm error, and its von Mises error \numHTwoBRel{} lower,
\numCmKnormVm{} against \numCmLabVm{} (guard \numHTwoBTau{}). On the seeds' medians the
guard gave the same readings, and the 95\% intervals of the two relative differences, from
Welch's statistic over the seeds and from resampling the instances, lay below zero. Instance
by instance it had the lower relative energy gap on \numPairCmKnormLabGap{} of the
instance-seed pairs, the lower von Mises error on \numPairCmKnormLabVm{} and the lower
displacement error on \numPairCmKnormLabDisp{}: changing the norm lowered the stress error
by more than half and left the displacement error about where it was, \numCmKnormDisp{}
against \numCmLabDisp{}. Weighted by element area, as in \cref{prop:stress}, instead of
equally, the ratio of the two networks' mean von Mises errors was \numSpecVmAreaRatio{},
against \numSpecVmElemRatio{} for the reported metric. On the held-out instances the
reduction of the von Mises error replicated, by more than half, and that of the mean
relative energy gap did not (\cref{sec:holdout}).

\subsection{Where the errors lie in the spectrum}\label{sec:norm:spectrum}

\Cref{rem:factor} writes the relative energy gap of a prediction as the product of the
squared relative displacement error $\delta^2$ and the normalised Rayleigh quotient $\rho$,
which measures how far up the spectrum of $K$ the error lies relative to the solution. After
the two-dimensional run, and without pre-registration, we measured both for the four
networks trained on 1,024 instances on each of the \numSpecLoads{} triples of seed,
validation instance and load case, and, from the eigendecomposition of each instance's
stiffness matrix, how each error's squared energy norm is distributed over the eigenmodes
(\cref{fig:spectra}). The solution itself lies at the bottom of the spectrum: its Rayleigh
quotient was a median \numSpecRqStarOverMin{} times the smallest eigenvalue of $K$. The
median of $\rho$ was \numSpecRqFree{} for the label-free transformer, \numSpecRqKnorm{} for
the energy-norm-supervised transformer, \numSpecRqLab{} for the displacement-supervised
transformer and \numSpecRqMgn{} for the graph network.

\begin{figure}[tbp]
\centering
\figfull{generated/fig_spectra.pdf}
\caption{Where the errors lie in the spectrum of the stiffness matrix (two dimensions,
networks trained on 1,024 instances; post hoc). Left: normalised Rayleigh quotient $\rho$ of
the error against the relative displacement error $\delta$, one point per seed, validation
instance and load case; large markers: medians; dashed lines: equal relative energy gap
$g=\rho\,\delta^2$. Right: median share of the error's squared energy norm in the
eigenmodes with $\lambda_m/R(\Ustar)$ at or above the abscissa ($\lambda_m$: eigenvalue;
$R(\Ustar)$: the solution's Rayleigh quotient); dashed curve: the solution's share.}
\label{fig:spectra}
\end{figure}

The two transformers that differ only in the norm of the loss had errors of about the same
size in different places. Over the triples, the geometric mean of the ratio of their $\rho$,
displacement-supervised over energy-norm-supervised, was \numSpecLabKnormRq{}
(\numSpecLabKnormRqSeeds{} in seeds 0, 1 and 2), and the displacement-supervised
transformer's $\rho$ was the larger on \numSpecLabKnormAbove{} of the triples; the geometric
mean of the ratio of their $\delta^2$ was \numSpecLabKnormDisp{}. On a logarithmic scale the
ratio of $\rho$ accounts for \numSpecSpectralShare{} of the geometric-mean ratio of their
relative energy gaps, \numSpecLabKnormGap{}, and the ratio of $\delta^2$ for the rest. The
modes with $\lambda_m\ge 100\,R(\Ustar)$ held a median \numSpecShareKLab{} of the squared
energy norm of the displacement-supervised transformer's error, against
\numSpecShareKKnorm{} for the energy-norm-supervised transformer, \numSpecShareKFree{} for
the label-free transformer and \numSpecShareKStar{} of the solution's own. The graph
network's errors had the size of the displacement-supervised transformer's (ratio of
$\delta^2$ \numSpecMgnLabDisp{}) and lay higher still (ratio of $\rho$ \numSpecMgnLabRq{}).
The networks ran with TF32 matrix products, whose rounding weighs heavily in the energy
norm; with TF32 off, the difference between the two norms reads larger still (supplementary
Section~\SuppRepro{}).

In three dimensions, where the full eigendecomposition was not computed, the same quotient
was measured after the run from the norms of the errors for the networks of the
three-dimensional comparison, not for the energy-norm-supervised transformer. Its median was
\numFieldRqFree{} for the label-free transformer, \numFieldRqLab{} for the
displacement-supervised transformer and \numFieldRqMgn{} for the graph network, and it was
larger for the displacement-supervised transformer than for the label-free one on
\numFieldRqAboveLab{} of the \numFieldLoads{} triples of seed, instance and load case.
Training in the Euclidean norm of the displacement thus left errors that lay higher in the
spectrum: of about the same size in two dimensions, and larger in three, where the
displacement-supervised transformer's displacement error also exceeded the label-free one's
(\cref{tab:3d}). By \eqref{eq:factor} such errors cost more energy and, by
\cref{prop:stress}, more stress error. Supplementary Section~\SuppTwo{} gives the
two-dimensional readings seed by seed, and the stress fields of one instance.
